\chapter*{Abstract}
\addcontentsline{toc}{chapter}{Abstract}

Air pollution is among the most severe environmental health threats facing Vietnam, yet only about forty regulatory stations across the country's 331,000 km\textsuperscript{2} supplied continuous PM2.5 over the study period --- thirty-seven after quality control --- leaving the majority of the population without local air-quality information. Satellite-augmented machine learning has been proposed as a remedy, with studies in densely monitored regions routinely reporting coefficients of determination above 0.8. This thesis asks whether such performance is attainable for Vietnam's sparse, tropical network, and it confronts a systematic methodological concern: the discrepancy between the random cross-validation used in most of the literature and spatial cross-validation, which reflects the operational task of predicting at unmonitored locations.

A 66-feature XGBoost-DART model was trained over 37 regulatory stations using MODIS and Himawari aerosol optical depth, TROPOMI trace gases, ERA5 meteorology, chemical-transport-model (CTM) outputs, and land-use predictors, and evaluated under leave-one-station-out (LOSO) validation and external validation against an independent low-cost-sensor (LCS) network and a reference-grade US Embassy monitor. The central result is a three-number arc: an R\textsuperscript{2} of approximately 0.80 under random cross-validation collapses to approximately 0.20 under honest LOSO, and to a deployable per-station median near 0.04 once oracle information is withdrawn. Tier stratification by pollution level (the T4F scheme) is the single largest performance lever, yielding the LOSO figure, but it carries a circular dependency --- assigning a station's tier requires knowing its mean PM2.5, the very quantity to be predicted --- and seven families of satellite-observable proxy methods failed to resolve it. The dependency is nonetheless resolved from the monitoring network rather than from observables: a spatial-prior routing model that anchors each station's baseline to the observed means of its neighbours recovers a deployable per-station mean R\textsuperscript{2} of 0.197 against an oracle ceiling of 0.203, with no dangerous low-to-high misclassification, conditional on proximity to the network. The CTM and existing global products fail distinctly: GEOS-CF shows an inverted diurnal cycle and a bias exceeding 200\%, MERRA-2 attains an Index of Agreement of only 0.39, and GHAP achieves roughly 50\% tier concordance. External validation is healthy within dense coverage, with an LCS median R\textsuperscript{2} of 0.53 and an Embassy R\textsuperscript{2} of 0.68, degrading with distance from the nearest anchor. The binding constraint is therefore station density rather than algorithmic capacity, and the densification of the network with calibrated low-cost sensors is the most tractable path to improved national accuracy.

\textbf{Keywords:} PM2.5, remote sensing, aerosol optical depth, machine learning, spatial cross-validation, Vietnam, low-cost sensors

